\chapter{Бозон Хиггса без тайн}

\lectureoverview{Механизм Хиггса лучше всего понимать в том порядке, в каком его составные части становятся необходимыми. Мы начнем с квантованного движения и полей, превратим обычную чашу в потенциал типа мексиканской шляпы и с помощью возникающего конденсата поймем, почему фермионы и слабые бозоны могут иметь массу. Лишь затем мы отождествим физический бозон Хиггса с радиальным колебанием этого конденсата и проследим процессы рождения и распада, которые привели к его открытию в 2012.}

\section{Механизм за заголовками}

Открытие новой частицы в области $125\,\mathrm{GeV}$ по праву вызвало большой интерес, но одного интереса еще недостаточно для понимания. Полезный вопрос состоит не в том, важен ли бозон Хиггса, и не в том, можно ли риторически связать его с происхождением Вселенной. Полезный вопрос гораздо уже: как устроено это явление в деталях и почему оно было нужно Стандартной модели?

Само название тоже требует краткой оговорки. Peter Higgs был одним из нескольких исследователей, выявивших соответствующую физику, а Fran\c{c}ois Englert и Robert Brout независимо пришли к центральной идее. Поэтому выражение \emph{Brout--Englert--Higgs mechanism} точнее отражает историю, хотя ниже будет использоваться более короткое название.

Объяснение разбито на модули. Сначала нам понадобятся два элементарных факта из квантовой механики и одна идея из классической теории поля. Затем нам понадобится вакуум, чье состояние наименьшей энергии содержит ненулевое поле. После этого мы сможем перечислить частицы Стандартной модели и увидеть, почему голая версия теории оставляет соответствующие частицы безмассовыми. Сам бозон Хиггса появляется только ближе к концу. Он не составляет весь механизм; это лишь одно конкретное колебание, остающееся после того, как механизм уже установлен.

\section{Два квантовых факта и поле}

Первый квантовый факт — это квантование. Момент импульса служит наглядным прототипом, потому что принимает дискретные значения,
\begin{equation}
L=n\hbar,
\qquad
n=0,\pm1,\pm2,\ldots .
\end{equation}
Встречаются и полуцелые значения момента импульса, но в ближайшем рассуждении от них ничего не зависит. Важно то, что непрерывно изменяющееся классическое движение превращается в лестницу квантовых состояний.

Второй факт — принцип неопределенности. Мы отложим его до того момента, когда появится непрерывная окружность возможных вакуумов. Тогда он не позволит нам рассматривать вакуум как маленький классический шарик, неподвижно лежащий в одной точно заданной точке.

Между этими двумя квантовыми фактами находится классическое понятие поля. Поле ставит в соответствие величину каждой точке пространства и времени. Электрические, магнитные и гравитационные поля — знакомые примеры. Область может выглядеть пустой и все же содержать поле. Представьте себе две заряженные пластины, расположенные так далеко, что ни одну из них не видно. Почти однородное электрическое поле между ними тем не менее локально вполне реально. Следовательно, ``пустое пространство'' не обязано означать ``каждое поле равно нулю.'' Это означает, что поля находятся в состоянии с наименьшей доступной энергией.

Предположим, что энергия поля минимальна при одном значении поля и возрастает, когда его смещают. Если поле возмутить, оно колеблется около минимума. Квантовая механика квантует эти колебания, и отдельные кванты являются частицами. На протяжении всей лекции нам понадобится следующий словарь теории поля:
\begin{equation}
\text{particle}
\quad\longleftrightarrow\quad
\text{quantum of a field vibration}.
\end{equation}

\section{Чаша в пространстве полей}

Для одной вещественной компоненты поля $\phi$ простейший энергетический ландшафт — это чаша,
\begin{equation}
V(\phi)=\frac{1}{2}m^{2}\phi^{2}.
\end{equation}
Дно чаши находится при $\phi=0$. Малые колебания около дна имеют определенную частоту, и их кванты ведут себя как частицы с массой $m$.

Теперь заменим одно поле двумя компонентами, $(\phi_1,\phi_2)$. Энергия может зависеть только от радиуса в этой внутренней плоскости,
\begin{equation}
r^{2}=\phi_1^{2}+\phi_2^{2},
\qquad
V=V(r).
\end{equation}
Это чаша над \emph{пространством полей}, а не углубление в обычном пространстве. Поворот пары $(\phi_1,\phi_2)$ не изменяет энергию. Шляпа, использованная в лекции, делает ту же симметрию наглядной: поверните ее вокруг оси, и она будет выглядеть по-прежнему.

\begin{figure}[tbp]
\centering
\includegraphics[width=0.95\textwidth]{lecture_144_figure_02.png}
\caption{Обычная симметричная чаша и настоящая шляпа, использованная для наглядной демонстрации вращательной симметрии в пространстве полей.}
\lectureframe{00:10:02}
\end{figure}

Аналогия становится чем-то большим, чем просто геометрией, когда поле движется. Сместите его из центра и пустите по кругу вдоль чаши. Это движение по кругу несет момент импульса в пространстве полей. Это не момент импульса объекта, вращающегося в лаборатории; это внутренний момент импульса. После квантования он дает модель сохраняющегося заряда,
\begin{equation}
Q=nq_0,
\qquad
n\in\mathbb{Z}.
\end{equation}
Уравнение схематично. Главная физическая идея состоит в том, что возбуждение заряженного поля можно представлять как вращение во внутренней плоскости. Положительный и отрицательный заряды соответствуют противоположным направлениям вращения, а отсутствию вращения соответствует нулевой заряд.

\section{Переворачиваем шляпу}

У обычной чаши есть только одна нижняя точка. Физика Хиггса начинается, когда шляпу переворачивают. Центр теперь неустойчив, тогда как весь круговой желоб имеет одну и ту же минимальную энергию. Стандартная математическая реализация имеет вид
\begin{equation}
V(\phi_1,\phi_2)
=
\lambda
\left(
  \phi_1^{2}+\phi_2^{2}-v^{2}
\right)^{2},
\qquad \lambda>0,
\end{equation}
а ее минимумы удовлетворяют условию
\begin{equation}
\phi_1^{2}+\phi_2^{2}=v^{2}.
\end{equation}
\editorialnote{Явный квартичный потенциал — это стандартное аналитическое завершение геометрии мексиканской шляпы, изображенной в лекции; он не переписан с формулы на доске.}

\begin{figure}[tbp]
\centering
\includegraphics[width=0.95\textwidth]{lecture_144_figure_03.png}
\caption{Центральный пик и круговой желоб потенциала перевернутой шляпы. Полный рисунок на доске заменяет более ранний обрезанный кадр, в котором отсутствовала большая часть геометрии.}
\lectureframe{00:13:10}
\end{figure}

Теперь вакуум уже не соответствует нулевому полю. Он должен лежать где-то на желобе,
\begin{equation}
(\phi_1,\phi_2)
=
v(\cos\theta,\sin\theta).
\end{equation}
Уравнения сохраняют вращательную симметрию, но любой конкретный вакуум выбирает один угол $\theta$. Это спонтанное нарушение симметрии: закон симметричен, а выбранное состояние наименьшей энергии — нет.

Теперь существуют два качественно разных движения. Радиальное движение внутрь или наружу ведет вверх по потенциалу и требует энергии. Движение вдоль желоба меняет только $\theta$ и на классическом уровне не требует потенциальной энергии. Поэтому с очень малой дополнительной потенциальной энергией может быть связано большое количество зарядоподобного углового движения.

\begin{figure}[tbp]
\centering
\includegraphics[width=0.95\textwidth]{lecture_144_figure_04.png}
\caption{Отмеченная точка в желобе представляет конкретный вакуум, выбранный из непрерывной окружности одинаково хороших минимумов.}
\lectureframe{00:16:14}
\end{figure}

\section{Неопределенность превращает вакуум в конденсат}

Классически можно было бы поместить поле в точку на ободе, положить его угловую скорость равной нулю и объявить задачу завершенной. Квантовая механика не допускает столь категоричного вывода. Координата и импульс подчиняются
\begin{equation}
\Delta x\,\Delta p\gtrsim\frac{\hbar}{2}.
\end{equation}
Аналогичное утверждение для кругового желоба состоит в том, что локализация угла поля делает неопределенным сопряженный ему внутренний момент импульса. В схематическом языке лекции,
\begin{equation}
\Delta\theta\,\Delta L_{\mathrm{int}}
\gtrsim
\frac{\hbar}{2}.
\end{equation}
\editorialnote{С операторами угла связаны технические тонкости, поэтому это соотношение — локальное полуклассическое выражение аргумента лекции о неопределенности, а не глобально точное операторное тождество.}

Поэтому вакуум, сосредоточенный вблизи одной точки желоба, не является состоянием с определенным внутренним зарядом. Он содержит широкую суперпозицию собственных состояний заряда. В лекции это уточняется операционным тестом: добавление одной единицы заряда или удаление одной единицы почти не меняет состояние. Схематически можно изобразить
\begin{align}
\ket{\Omega}
&\sim
\sum_n c_n\ket{n},
\\
c_{n+1}
&\simeq
c_n
\quad\text{within the occupied range}.
\end{align}
Тогда сдвиг $n$ на единицу по существу лишь переобозначает члены, уже присутствующие. В этом смысле вакуум является \emph{конденсатом}: он содержит неопределенное число соответствующих квантов, так что добавление или удаление одного из них не дает наблюдаемо иного макроскопического состояния.

Реальный электрослабый вакуум не является конденсатом обычного электрического заряда. Если бы это было так, электромагнитные свойства пустого пространства радикально отличались бы. Однако сверхпроводник дает конкретный пример заряженного конденсата. Поэтому он служит полезным физическим мостом между абстрактной картиной мексиканской шляпы и утверждением, что вакуум может обладать когерентной квантовой структурой.

Этим завершается первый модуль. Мы узнали, что состояние наименьшей энергии не обязано находиться при нулевом поле и что квантовая неопределенность может сделать это состояние нечувствительным к добавлению или удалению отдельных квантов. Теперь спросим, каким частицам нужен такой фон.

\section{Голая Стандартная модель}

Естественным верхним ориентиром для масс элементарных частиц является планковская масса. Если достаточно сильно сконцентрировать локализованную энергию, гравитацией уже нельзя будет пренебречь; в конце концов объект правильнее рассматривать как черную дыру, а не как обычную элементарную частицу. Численно,
\begin{equation}
m_{\mathrm P}
=
\sqrt{\frac{\hbar c}{G}}
\approx
2.2\times10^{-8}\,\mathrm{kg},
\end{equation}
что примерно соответствует массе крошечной пылинки. Известные элементарные частицы фантастически легче. В масштабной оценке, использованной в лекции, их массы достигают лишь примерно $10^{-17}$ планковской массы.

Это наблюдение не доказывает, что более тяжёлых элементарных частиц не существует. Ускорители преимущественно обнаруживают частицы, достаточно лёгкие, чтобы их можно было произвести при имеющейся энергии. Поэтому известный спектр отчасти определяется досягаемостью установок, использованных для его поиска. За пределами нынешней досягаемости может существовать много более тяжёлых частиц.

Наблюдаемые частицы делятся на фермионы и бозоны. К фермионам относятся электрон, мюон, тау, соответствующие им нейтрино и шесть кварковых ароматов. К бозонам относятся фотон, глюоны и слабые бозоны $W$ и $Z$. Глюоны связывают кварки в протоны и нейтроны и потому играют центральную роль в структуре обычных ядер.

\begin{figure}[tbp]
\centering
\includegraphics[width=0.95\textwidth]{lecture_144_figure_05.png}
\caption{Первая в лекции сводка Стандартной модели, с огромным диапазоном масс, показанным над столбцами фермионов и бозонов.}
\lectureframe{00:24:22}
\end{figure}

Теория — не просто список; это набор разрешённых взаимодействий. Электрон может испустить фотон,
\begin{equation}
e\longrightarrow e+\gamma.
\end{equation}
Кварк несёт и электрический заряд, и цветовой заряд, поэтому он может испустить либо фотон, либо глюон,
\begin{equation}
q\longrightarrow q+\gamma,
\qquad
q\longrightarrow q+g.
\end{equation}
Электрон не может испустить глюон, а нейтрино не испускает ни фотон, ни глюон. Фермионы с подходящими слабыми квантовыми числами могут взаимодействовать через $Z$,
\begin{equation}
f\longrightarrow f+Z.
\end{equation}
Бозоны $W$ и реакции с заряженным током тоже входят в Стандартную модель, но в лекции они опущены, чтобы всю аргументацию можно было провести с помощью одного нейтрального слабого бозона.

\begin{figure}[tbp]
\centering
\includegraphics[width=0.95\textwidth]{lecture_144_figure_06.png}
\caption{Вершины испускания фотона, глюона и $Z$ рядом со словом ``безмассовые.'' Эти процессы резюмируют голую теорию до введения конденсата.}
\lectureframe{00:29:20}
\end{figure}

Фотон безмассов. В ненарушенной калибровочной теории того же общего типа глюоны и слабые калибровочные бозоны тоже безмассовы. Немодифицированная электрослабая теория точно так же не может просто ввести массы фермионов вручную, сохраняя свою калибровочную симметрию. Не хватает одного ингредиента. Прежде чем ввести его, нам нужно увидеть в более простой обстановке, как фоновое поле может изменить массу частицы, не действуя подобно трению.

\section{Как фоновое поле может сдвигать массу}

Рассмотрим молекулу воды как электрический диполь: один конец слегка положителен, а другой слегка отрицателен. В пустом пространстве две противоположные ориентации энергетически эквивалентны. В лекции этим ориентациям шутливо дают разные названия, ``вода'' и ``виски'', только для того, чтобы за ними можно было следить как за различными внутренними состояниями.

Поместим молекулу между пластинами конденсатора. Если её дипольный момент равен $\bm d$, а электрическое поле равно $\bm E$, то энергия взаимодействия равна
\begin{equation}
U=-\bm d\mathbin{\cdot}\bm E.
\end{equation}
Для двух ориентаций, направленных вдоль и против поля, это принимает вид
\begin{equation}
U_{\parallel}=-dE,
\qquad
U_{\mathrm{anti}}=+dE.
\end{equation}
Соотношение Эйнштейна $E=mc^2$ превращает расщепление по энергии в расщепление по массе,
\begin{equation}
m_{\mathrm{anti}}-m_{\parallel}
=
\frac{2dE}{c^{2}}.
\end{equation}
О молекулу ничто не трётся. Просто одно внутреннее состояние имеет большую энергию и, следовательно, большую массу, чем другое, пока присутствует фоновое поле.

\begin{figure}[tbp]
\centering
\includegraphics[width=0.95\textwidth]{lecture_144_figure_07.png}
\caption{Противоположные ориентации молекулы в поле конденсатора. Поле расщепляет их энергии, не создавая силы сопротивления.}
\lectureframe{00:34:35}
\end{figure}

Это устраняет две вводящие в заблуждение аналогии. Поле Хиггса — не патока, сквозь которую частицы должны продираться, и частица с массой — не снегоочиститель, замедляемый веществом, скапливающимся перед ним. В пространственно однородном поле вообще может не быть никакой результирующей силы. Эффект состоит во внутреннем сдвиге энергии, а не в трении и не в потере импульса.

Тот же пример можно пересказать на языке квантовой теории. Классическое электрическое поле — это когерентное состояние, которое для наших целей можно описать как конденсат фотонов. Молекула может испустить фотон в этот фон или поглотить из него фотон. Поскольку добавление или удаление одного кванта почти не изменяет когерентное состояние с большим числом квантов, линия фотона может заканчиваться на конденсате. Крест, нарисованный на доске, обозначает это исчезновение в фоне, а не нарушение закона сохранения энергии.

\begin{figure}[tbp]
\centering
\includegraphics[width=0.95\textwidth]{lecture_144_figure_08.png}
\caption{Поле конденсатора, перерисованное как фотонный конденсат. Фотон, испущенный диполем, поглощается фоном, нечувствительным к одному дополнительному кванту.}
\lectureframe{00:38:45}
\end{figure}

Это именно та схема, которая нам нужна. Если конденсат может поглотить разность квантовых чисел между двумя внутренними состояниями, он может смешивать эти состояния. Возникающее при этом осциллирование может проявляться как масса.

\section{Почему Хиггс не объясняет всю массу}

Прежде чем применять эту схему, нужно устранить одно заблуждение. Неверно, что всякий массивный объект получает всю свою массу от поля Хиггса.

Поместите излучение в идеально отражающий ящик. Отдельные фотоны безмассовы, но для замкнутого ящика существует система покоя, и в нём запасена энергия. Его масса увеличивается на
\begin{equation}
\Delta M=\frac{E_{\mathrm{radiation}}}{c^{2}}.
\end{equation}
Следовательно, масса может возникать из заключённой кинетической и полевой энергии, даже когда составляющие не имеют массы покоя.

Важный реальный пример — протон. Он содержит кварки, антикварки и глюоны в состоянии сильного взаимодействия. Большая часть массы протона возникает из их кинетической энергии и энергии связи, а не из простого сложения хиггсовских масс нескольких кварков-составляющих. В оценке, подчёркнутой в лекции, если сделать лёгкие кварки безмассовыми, масса протона изменилась бы лишь на уровне порядка процента. Чёрные дыры ещё резче подчёркивают логическую суть: чёрная дыра может быть массивной из-за полной энергии, из которой она образовалась, а не потому, что каждый вклад сначала был превращён полем Хиггса в массу покоя элементарных частиц.

У механизма Хиггса есть более точная задача. Он допускает массы элементарных фермионов и слабых калибровочных бозонов, сохраняя при этом непротиворечивость лежащей в основе калибровочной теории.

\section{Хиральность: масса как смешивание левого--правого}

Сосредоточимся сначала на электроне. Спин — это собственный момент импульса. Для движущейся частицы можно сравнить направление спина с направлением импульса. В геометрическом языке лекции правый электрон и левый электрон различаются соотношением между этими направлениями. Для безмассовой частицы спиральность и хиральность соотносятся тем простым образом, который нужен для рассуждения; для массивной частицы эти точные понятия необходимо различать.

Представьте, что вы снова и снова обгоняете массивный электрон. В другой инерциальной системе отсчёта его импульс меняет направление на противоположное, а спин — нет, так что кажущаяся рукость, связанная с направлением движения, может обратиться. Ни один наблюдатель не может обогнать строго безмассовую частицу, которая всегда движется со скоростью света. Это мотивирует эвристическое утверждение лекции: безмассовый электрон остаётся в одном хиральном состоянии, тогда как массивный электрон непрерывно смешивает левую и правую хиральности.

Точное утверждение теории поля состоит в том, что массовый член Дирака связывает два хиральных поля,
\begin{equation}
\mathcal L_m
=
-m_e\bar e e
=
-m_e
\left(
  \bar e_L e_R
  +
  \bar e_R e_L
\right).
\end{equation}
\editorialnote{Член массы Дирака — это стандартное математическое завершение лекционной картины ``частоты переворота''. Хиральное смешивание пропорционально массе; буквальная частота таких переворотов зависит от состояния и от релятивистского замедления времени.}

Если $m_e=0$, левое и правое поля распространяются независимо. Если $m_e\ne0$, уравнения связывают их. В этом точном смысле масса измеряет силу смешивания левого--правого.

\section{Закон сохранения, запрещающий голую массу}

Почему бы просто не записать член смешивания с самого начала? Слабое взаимодействие по-разному обращается с $e_L$ и $e_R$. Сасскинд намеренно даёт этому различию простое название: \emph{зилч}. В нормировке, использованной в лекции,
\begin{equation}
z(e_L)=1,
\qquad
z(e_R)=0.
\end{equation}
У двух хиральностей одинаковый электрический заряд, но разные слабые квантовые числа. Поэтому прямой переход
\begin{equation}
e_L\longrightarrow e_R
\end{equation}
не сохранял бы величину, различаемую слабым взаимодействием. В ненарушенной теории голый массовый член запрещён.

\editorialnote{``Зилч'' — педагогическое сокращение. В Стандартной модели левые и правые фермионы принадлежат разным представлениям $SU(2)_L\times U(1)_Y$, поэтому прямой массовый член Дирака не является калибровочно-инвариантным. Лекционные присваивания $1$ и $0$ кодируют это различие, не воспроизводя полной таблицы электрослабых зарядов.}

На доске это препятствие представлено как цепочка. Правый электрон не может стать левым, если только что-то ещё не унесёт разность зилча. То же препятствие относится к кваркам, мюонам и тау, причём после введения недостающего поля сила такого перехода зависит от вида частицы.

\begin{figure}[tbp]
\centering
\includegraphics[width=0.95\textwidth]{lecture_144_figure_09.png}
\caption{Хиральная цепочка на доске: квант конденсата, вставленный в зелёном кресте, позволяет метке зилча измениться при переходе между право- и левохиральным распространением.}
\lectureframe{00:50:15}
\end{figure}

\section{Конденсат ``Зиггса'' и масса фермионов}

Введём поле, кванты которого несут в точности недостающее слабое квантовое число. Чтобы не возникло преждевременной путаницы с физической частицей Хиггса, в лекции соответствующий квант называется \emph{зиггс}. Предположим, что поле зиггса имеет потенциал типа «мексиканская шляпа», а его вакуум представляет собой конденсат. Тогда левохиральный электрон может испустить квант зиггса и стать правохиральным,
\begin{equation}
e_L\longrightarrow e_R+\text{Ziggs},
\end{equation}
или правохиральный электрон может поглотить такой квант и стать левохиральным,
\begin{equation}
e_R+\text{Ziggs}\longrightarrow e_L.
\end{equation}
Конденсат поглощает эту разность квантовых чисел или восполняет её, не меняя своего макроскопического состояния. То, что было запрещено в пустом вакууме с ненарушенной симметрией, становится разрешённым в вакууме с нарушенной симметрией.

Повторные испускания и поглощения смешивают $e_L$ и $e_R$. Поэтому электрон распространяется как массивная частица. Тот же механизм действует для заряженных лептонов и кварков. У каждого вида своя связь с конденсатом, так что сила переворота и, следовательно, масса могут различаться от вида к виду.

В современной записи калибровочно-инвариантное взаимодействие представляет собой связь Юкавы. После того как поле Хиггса приобретает вакуумное среднее значение $v$, получается
\begin{equation}
\mathcal L_Y
\supset
-\frac{y_f}{\sqrt{2}}
(v+h)\bar f f,
\end{equation}
и, следовательно,
\begin{equation}
m_f=\frac{y_fv}{\sqrt{2}},
\qquad
g_{hff}=\frac{y_f}{\sqrt{2}}
=
\frac{m_f}{v}.
\end{equation}
\editorialnote{Эти формулы дают стандартную электрослабую запись, стоящую за лекционным рассуждением о связи с конденсатом. Они также явно показывают, почему одна и та же константа связи определяет и массу фермиона, и распад Хиггса на этот фермион.}

В лекции это описывается как спонтанное нарушение хиральной симметрии. Это передаёт лево--правую картину на уровне примера. Точнее, вакуумное среднее поля Хиггса нарушает электрослабую калибровочную симметрию до электромагнитной подгруппы, а взаимодействия Юкавы затем становятся массовыми членами фермионов.

\section{Как слабый бозон приобретает массу}

Конденсат изменяет и слабые калибровочные бозоны. Возбуждение зиггса может испустить бозон $Z$, и наоборот, $Z$ может переходить в соответствующее возбуждение, поддерживаемое конденсатом, и обратно. Тем самым $Z$ распространяется в вакууме, где он смешивается с угловым сектором поля Хиггса. Результатом становится массивный векторный бозон с продольной поляризацией, которой безмассовый векторный бозон не имел бы.

\begin{figure}[tbp]
\centering
\includegraphics[width=0.95\textwidth]{lecture_144_figure_10.png}
\caption{Завершённая последовательность на доске для механизма слабого бозона: хиральная цепочка остаётся сверху, а $Z$ внизу обменивается квантовым числом с сектором, поддерживаемым конденсатом.}
\lectureframe{00:53:40}
\end{figure}

Это явление Броута--Энглерта--Хиггса в его калибровочно-бозонной форме. В полной Стандартной модели массу приобретают как бозоны $W$, так и $Z$, тогда как одна комбинация исходных нейтральных калибровочных полей остаётся безмассовым фотоном. Если бы вакуум вместо этого нарушал обычную электромагнитную калибровочную симметрию, фотон стал бы массивным, как это фактически происходит внутри сверхпроводника. Тогда у пустого пространства было бы совершенно иное дальнодействующее электромагнитное поведение.

Слабые бозоны уже задолго до выделения физической частицы Хиггса дали веские свидетельства в пользу этого механизма. Их массы и взаимодействия показали, что электрослабая симметрия была нарушена требуемым образом. Поэтому в лекции можно сказать, в некотором свободном операциональном смысле, что сектор ``зиггс'' уже был известен: недостающим экспериментальным объектом оставалось самостоятельное радиальное колебание.

\editorialnote{Лекционное импровизированное упоминание открытия в SLAC около 1980 года сжимает в одно несколько событий. Нейтральные слабые токи были обнаружены в 1973 году, нарушающее чётность рассеяние электронов в SLAC было зарегистрировано в конце 1970-х, а бозоны $W$ и $Z$ были непосредственно открыты в CERN в 1983 году. Физически существенен здесь более ранний экспериментальный факт нарушения электрослабой симметрии.}

\section{Чем на самом деле является бозон Хиггса}

Вернёмся к «мексиканской шляпе». Движение вдоль желоба меняет угол, но не радиус. Именно этот угловой сектор даёт моды, поглощаемые массивными слабыми бозонами. Физический бозон Хиггса соответствует другому движению: колебанию внутрь и наружу, вверх по стенкам потенциала.

Схематически запишем поле как
\begin{equation}
\Phi(x)
=
\frac{v+h(x)}{\sqrt{2}}
\exp\!\left(
  \frac{i\theta(x)}{v}
\right).
\end{equation}
Переменная $\theta$ описывает движение вдоль желоба. Переменная $h$ меняет радиус. В калибровочной теории соответствующие угловые степени свободы становятся продольными компонентами массивных слабых бозонов, тогда как $h$ остаётся распространяющейся скалярной частицей: бозоном Хиггса.

\begin{figure}[tbp]
\centering
\includegraphics[width=0.95\textwidth]{lecture_144_figure_11.png}
\caption{На доске вновь появляются оба потенциала, причём в желобе «мексиканской шляпы» отмечено радиальное смещение. Это колебание внутрь и наружу и есть физическая мода Хиггса.}
\lectureframe{00:56:30}
\end{figure}

Существует и эквивалентная многочастичная картина. Вакуум представляет собой конденсат, и радиальное колебание изменяет плотность конденсата. Поэтому бозон Хиггса подобен волне сжатия в этой квантовой среде. Аналогия полезна, если только не делать из неё вывода, будто такая волна обязана быть стабильной. Обычные звуковые волны могут затухать, и возбуждение Хиггса тоже может распадаться.

Это различие объясняет, почему обнаружение бозона Хиггса стало последним недостающим звеном, а не первым свидетельством механизма. Вакуум с нарушенной симметрией и массивные слабые бозоны уже были косвенно видны. Радиальную волну плотности нужно было породить как частицу и опознать по продуктам её распада.

\section{Распад, рождение и преимущество тяжёлых фермионов}

Хиггс взаимодействует с фермионом с силой, пропорциональной массе этого фермиона. Он может распадаться на пару фермион-антифермион,
\begin{equation}
H\longrightarrow f+\bar f,
\end{equation}
и ту же диаграмму можно читать в обратном направлении как рождение,
\begin{equation}
f+\bar f\longrightarrow H.
\end{equation}
Поэтому рождение Хиггса в аннигиляции электрона и позитрона в принципе возможно. Исторически, однако, у ранних коллайдеров сначала просто не хватало энергии для рождения Хиггса, а более поздние машины столкнулись с более фундаментальной проблемой: юкавская связь электрона чрезвычайно мала. Аннигиляция лёгких кварков столь же слаба.

\begin{figure}[tbp]
\centering
\includegraphics[width=0.95\textwidth]{lecture_144_figure_12.png}
\caption{Изображённая на доске вершина распада Хиггса. Если читать стрелки в обратную сторону, получается слияние фермиона и антифермиона в Хиггс.}
\lectureframe{00:59:40}
\end{figure}

Чем тяжелее фермион, тем больше его связь Юкавы. Это выделяет топ-кварк, масса которого составляет около $173\,\mathrm{GeV}$, то есть примерно $180$ протонных масс. В лекции используется округлённая оценка — около $170$ протонных масс. Однако Хиггс с массой около $125\,\mathrm{GeV}$ слишком лёгок, чтобы распасться на пару $t\bar t$ на массовой оболочке,
\begin{equation}
m_H<2m_t.
\end{equation}
Нельзя и построить пригодный пучок топ-кварков: топ распадается раньше, чем успевает адронизоваться в долгоживущую частицу пучка.

Квантовая механика даёт способ обойти оба препятствия. Два глюона могут через виртуальную петлю топ-кварка перейти в бозон Хиггса,
\begin{equation}
g+g\longrightarrow H.
\end{equation}
Виртуальные топовые линии не обязаны удовлетворять порогу на массовой оболочке, который действовал бы для рождения реальной топ-пары. Протоны полны глюонов, так что высокоэнергетические протон-протонные столкновения на Большом адронном коллайдере дают интенсивный источник требуемых пар глюонов.

\begin{figure}[tbp]
\centering
\includegraphics[width=0.95\textwidth]{lecture_144_figure_13.png}
\caption{Слияние глюонов через виртуальную топовую петлю. Два глюона, взятые из сталкивающихся протонов, встречаются в петле и рождают Хиггс.}
\lectureframe{01:06:35}
\end{figure}

Вот почему протонный коллайдер, несмотря на сложность протонных остатков, стал решающей фабрикой Хиггса. Протон не взаимодействует с Хиггсом как единый элементарный объект; его глюонные составляющие рождают Хиггс через петлю, в которой доминирует юкавская связь тяжёлого топ-кварка.

\section{Частица 2012 года и её проверка по двухфотонному каналу}

Частица, о которой сообщили в 2012 году, имела массу около
\begin{equation}
m_H\simeq125\,\mathrm{GeV}.
\end{equation}
В удобном для лекции сравнении это примерно $127$ протонных масс. Её открытие завершило список частиц, требуемых минимальной Стандартной моделью. Оставался вопрос, совпадают ли её количественные свойства в точности со свойствами Стандартной модели.

Логику топовой петли можно обратить. Хиггс связан с виртуальными заряженными частицами в петле, и петля может испустить два фотона,
\begin{equation}
H\longrightarrow\gamma+\gamma.
\end{equation}
Этот распад редок, но фотоны в адронном коллайдере являются исключительно чистыми объектами. Поэтому узкий избыток в распределении по их инвариантной массе может выявить Хиггс даже тогда, когда более частые распады экспериментально менее чисты.

\begin{figure}[tbp]
\centering
\includegraphics[width=0.95\textwidth]{lecture_144_figure_14.png}
\caption{Двухфотонный треугольник справа на доске, сопоставленный с более ранней петлёй глюонного слияния слева. Эти два процесса проверяют одну и ту же петлевую логику в противоположных направлениях.}
\lectureframe{01:09:00}
\end{figure}

Во время этой лекции сообщавшаяся частота двухфотонных событий казалась, возможно, примерно в полтора раза больше простейшего ожидания Стандартной модели, при значимости лишь около двух стандартных отклонений. Само по себе такое флуктуационное отклонение не было свидетельством. Однако если бы оно сохранилось, то могло бы указывать на дополнительную заряженную частицу в петле, возможно связанную с суперсимметрией или с другим расширением Стандартной модели.

Это утверждение относится к историческому моменту 2012. Оно фиксирует то, за чем экспериментаторам тогда говорили следить дальше, а не нынешнее заявление об аномалии. Само открытие в значительной степени подтвердило схему, уже подкреплённую данными о слабых бозонах. Более точными в научном отношении проверками были скорости распада: взаимодействовал ли новый скаляр с каждой частицей по зависящей от массы схеме, предсказанной механизмом Хиггса?

\section{Вопросы из аудитории}

\begin{classroomqa}
\audiencequestion{Чем определяется, почему один фермион осциллирует между левым и правым состояниями быстрее другого и потому имеет другую массу?\lecturetimestamp{01:12:00}}
\lecturerresponse{У каждого фермионного вида есть собственная константа связи с конденсатом Хиггса. Более сильная связь приводит к более сильному лево--правому смешиванию и большей массе. Та же константа связи определяет, насколько интенсивен распад физического бозона Хиггса в канал с данным фермионом, поэтому более тяжёлые фермионы, как правило, сильнее связаны с ним. Стандартная модель параметризует эти юкавские константы связи, но не объясняет их поразительную иерархию. Сказать, что у электрона малая константа связи, а у топ-кварка большая, — это точное описание, но пока не более глубокое объяснение.}
\end{classroomqa}

\begin{classroomqa}
\audiencequestion{Знаем ли мы величину поля Хиггса в вакууме?\lecturetimestamp{01:13:36}}
\lecturerresponse{Да. В используемой здесь конвенции вакуумное среднее значение составляет около $240\,\mathrm{GeV}$; более точное стандартное значение равно приблизительно $246\,\mathrm{GeV}$. Его можно представлять либо как радиальное смещение поля от начала координат, либо как плотность конденсата. Физический квант Хиггса — это малое радиальное колебание этой величины, или, что эквивалентно, локальная флуктуация плотности конденсата.}
\end{classroomqa}

\begin{classroomqa}
\audiencequestion{Означает ли ненулевое вакуумное значение поля Хиггса наличие плотности энергии пустого пространства?\lecturetimestamp{01:14:12}}
\lecturerresponse{Да. Потенциал Хиггса вносит вклад в вакуумную энергию, как и многие другие секторы квантовой теории поля. Наблюдаемая суммарная вакуумная энергия чрезвычайно мала по сравнению с естественными масштабами этих вкладов, так что они должны почти компенсировать друг друга. Ни одна общепринятая теория не объясняет этой компенсации. Это проблема космологической постоянной, одна из самых глубоких нерешённых проблем фундаментальной физики.}
\end{classroomqa}

\section{Чего удалось добиться этой конструкцией}

Теперь эту логику можно проследить от начала до конца, не прибегая к липкому эфиру. Квантованное внутреннее вращение служит моделью заряда. Превращение чаши в мексиканскую шляпу придаёт полю ненулевое вакуумное значение. Принцип неопределённости превращает точно выбранный классический вакуум в конденсат, нечувствительный к отдельным квантам. Тогда конденсат может нести различие по слабому квантовому числу между левохиральными и правохиральными фермионами, делая возможными их смешивание и тем самым возникновение масс. Он также обеспечивает продольные моды, которые позволяют $W$ и $Z$ распространяться как массивные векторные бозоны, тогда как фотон остаётся безмассовым.

Физический бозон Хиггса — это радиальное колебание, остающееся после этой перестройки. Это не ``сама масса'', и он не отвечает за бо́льшую часть массы протона. Это наблюдаемая плотностная волна поля, нарушающего симметрию. Его константы связи несут память о массах, порождённых этим полем, и именно поэтому в его рождении и распадах доминируют тяжёлые частицы. Поэтому петля глюонного слияния, дифотонный канал и последний вопрос из аудитории о вакуумной энергии — не отдельные курьёзы. Это следствия, проверки и ещё не завершённые выводы из одного и того же нетривиального вакуума.
